%=============================================================
% Capitolo 8 — Autovalori e Decomposizione SVD
%=============================================================
\chapter{Autovalori e Decomposizione ai Valori Singolari}
\label{ch:autovalori}

\section{Autovalori e autovettori}
\label{sec:autovalori}

\begin{definition}
Sia $A \in \mathbb{C}^{n \times n}$. Un numero $\lambda \in \mathbb{C}$ è un
\emph{autovalore} di $A$ e $v \in \mathbb{C}^n \setminus \{0\}$ il corrispondente
\emph{autovettore} se
\begin{equation}
  Av = \lambda v.
  \label{eq:autovettore}
\end{equation}
\end{definition}

\subsection{Equazione caratteristica}

$\lambda$ è autovalore di $A$ se e solo se $\det(A - \lambda I) = 0$.
Il polinomio $p(\lambda) = \det(A - \lambda I)$ di grado $n$ è detto
\textbf{polinomio caratteristico}.

\begin{remark}
In generale, il calcolo degli autovalori tramite il polinomio caratteristico è
\emph{numericamente instabile} per $n \geq 5$ (teorema di Abel-Ruffini). In pratica
si utilizzano metodi iterativi (ad esempio l'algoritmo QR).
\end{remark}

\subsection{Proprietà fondamentali}

Per $A \in \mathbb{C}^{n \times n}$ con autovalori $\lambda_1, \ldots, \lambda_n$:
\begin{align}
  \det(A)  &= \prod_{i=1}^n \lambda_i, \label{eq:det_autoval}\\
  \operatorname{tr}(A) &= \sum_{i=1}^n \lambda_i. \label{eq:traccia}
\end{align}

\subsection{Molteplicità algebrica e geometrica}

\begin{definition}
\begin{itemize}
  \item La \emph{molteplicità algebrica} $\mu_a(\lambda)$ è la molteplicità di
        $\lambda$ come radice del polinomio caratteristico.
  \item La \emph{molteplicità geometrica} $\mu_g(\lambda) = \dim\ker(A-\lambda I)$
        è il numero di autovettori linearmente indipendenti.
\end{itemize}
Si ha sempre $\mu_g(\lambda) \leq \mu_a(\lambda)$.
\end{definition}

%------------------------------------------------------------
\section{Raggio spettrale e cerchi di Gershgorin}
\label{sec:gershgorin}
%------------------------------------------------------------

\begin{theorem}[Cerchi di Gershgorin]
\label{thm:gershgorin}
Ogni autovalore di $A$ appartiene all'unione dei \emph{cerchi di Gershgorin}:
\begin{equation}
  \bigcup_{i=1}^n \mathcal{D}_i, \quad \text{con } \mathcal{D}_i =
  \left\{ z \in \mathbb{C} : |z - a_{ii}| \leq r_i \right\},\quad
  r_i = \sum_{j \neq i} |a_{ij}|.
  \label{eq:gershgorin}
\end{equation}
\end{theorem}

I cerchi di Gershgorin forniscono una \emph{stima rapida e grossolana} della
localizzazione degli autovalori, utile come sanity check prima di calcoli più
accurati.

%------------------------------------------------------------
\section{Diagonalizzazione}
\label{sec:diagonalizzazione}
%------------------------------------------------------------

\begin{definition}[Matrice diagonalizzabile]
$A \in \mathbb{C}^{n \times n}$ è \emph{diagonalizzabile} se esiste una matrice
invertibile $V$ tale che
\begin{equation}
  A = V \Lambda V^{-1}, \quad \Lambda = \operatorname{diag}(\lambda_1, \ldots, \lambda_n).
  \label{eq:diag}
\end{equation}
Le colonne di $V$ sono gli autovettori di $A$.
\end{definition}

\begin{theorem}[Matrici hermitiane]
Se $A = A^*$ (hermitiana; reale: $A = A^\top$), allora:
\begin{enumerate}
  \item Tutti gli autovalori sono \emph{reali};
  \item Gli autovettori relativi ad autovalori distinti sono \emph{ortogonali};
  \item $A$ è \emph{unitariamente diagonalizzabile}: $A = Q\Lambda Q^*$ con $Q$
        unitaria.
\end{enumerate}
\end{theorem}

%------------------------------------------------------------
\section{Decomposizione ai valori singolari (SVD)}
\label{sec:svd}
%------------------------------------------------------------

La SVD estende la diagonalizzazione a matrici \emph{rettangolari}.

\begin{theorem}[Decomposizione ai valori singolari]
\label{thm:svd}
Per ogni $A \in \R^{m \times n}$ esiste la fattorizzazione
\begin{equation}
  A = U \Sigma V^\top,
  \label{eq:svd}
\end{equation}
dove:
\begin{itemize}
  \item $U \in \R^{m \times m}$ è ortogonale (le colonne $u_i$ sono i
        \emph{vettori singolari sinistri});
  \item $V \in \R^{n \times n}$ è ortogonale (le colonne $v_i$ sono i
        \emph{vettori singolari destri});
  \item $\Sigma = \operatorname{diag}(\sigma_1, \ldots, \sigma_p) \in \R^{m \times n}$
        con $p = \min(m,n)$ e $\sigma_1 \geq \sigma_2 \geq \cdots \geq \sigma_p \geq 0$
        (\emph{valori singolari}).
\end{itemize}
\end{theorem}

\subsection{Relazione tra SVD e autovalori}

I valori singolari di $A$ sono le radici quadrate degli autovalori non negativi
di $A^\top A$ (e di $AA^\top$):
\begin{equation}
  \sigma_i = \sqrt{\lambda_i(A^\top A)} = \sqrt{\lambda_i(AA^\top)}.
  \label{eq:svd_autoval}
\end{equation}
Se $A$ è simmetrica definita positiva, $\sigma_i = \lambda_i$.

\subsection{SVD troncata e approssimazione di rango basso}

\begin{theorem}[Teorema di Eckart-Young]
La migliore approssimazione di rango $k$ di $A$ nella norma spettrale e nella
norma di Frobenius è:
\begin{equation}
  A_k = \sum_{i=1}^k \sigma_i\, u_i v_i^\top,
  \label{eq:svd_troncata}
\end{equation}
con errore $\|A - A_k\|_2 = \sigma_{k+1}$.
\end{theorem}

\subsection{Applicazioni della SVD}

\begin{itemize}
  \item \textbf{Norma spettrale}: $\|A\|_2 = \sigma_1$.
  \item \textbf{Numero di condizionamento}: $\kappa_2(A) = \sigma_1 / \sigma_n$ (per
        $A$ quadrata non singolare).
  \item \textbf{Rango numerico}: $\operatorname{rank}(A) = $ numero di valori singolari
        significativamente non nulli.
  \item \textbf{Pseudoinversa di Moore-Penrose}: $A^+ = V\Sigma^+ U^\top$ dove
        $\Sigma^+$ inverte i valori singolari non nulli.
  \item \textbf{Compressione dati}: la SVD troncata è usata in PCA, compressione di
        immagini, sistemi di raccomandazione.
\end{itemize}

%------------------------------------------------------------
\section{Algoritmo QR per il calcolo degli autovalori}
\label{sec:algo_qr}
%------------------------------------------------------------

L'\textbf{algoritmo QR} è il metodo standard per calcolare tutti gli autovalori
di una matrice:

\begin{algorithm}[htbp]
  \caption{Algoritmo QR (forma base)}
  \label{alg:qr_autoval}
  \begin{algorithmic}[1]
    \Require $A_0 = A \in \R^{n \times n}$, numero iterazioni $N$
    \Ensure Approssimazione degli autovalori di $A$
    \For{$k = 0, 1, 2, \ldots, N$}
      \State Calcola la fattorizzazione $A_k = Q_k R_k$
      \State $A_{k+1} \gets R_k Q_k$
    \EndFor
    \State \Return gli elementi diagonali di $A_N$
  \end{algorithmic}
\end{algorithm}

La sequenza $A_k$ converge (sotto opportune ipotesi) a una matrice triangolare
superiore (o quasi-triangolare reale) con gli autovalori sulla diagonale.

In pratica, prima si riduce $A$ a forma di Hessenberg superiore
(costo $\mathcal{O}(n^3)$), poi si applica l'algoritmo QR alla forma ridotta
(ogni iterazione costa $\mathcal{O}(n^2)$).
